% !TEX root = ../../main.tex
\subsection{系统参与方与信任边界}

FoodShield 系统涉及用户端、骑手端、平台服务器、管理员端和审计日志系统五类主要参与方。
不同参与方具有不同的业务权限和信息可见范围。
系统通过权限隔离、字段隐藏和数据最小化原则限制敏感信息暴露，避免普通通信功能演变为身份泄露通道。

\begin{table}[H]
\centering
\caption{系统主要参与方与权限边界}
\renewcommand{\arraystretch}{1.55}
\setlength{\tabcolsep}{5pt}
\zihao{-4}

\begin{tabular}{L{2.2cm} L{4.7cm} L{6.7cm}}
\toprule
\textbf{参与方} & \textbf{主要功能} & \textbf{权限与隐私边界} \\
\midrule

用户端
&
注册登录；创建订单；查看订单历史；订单搜索；添加备注标签；发起投诉申请。
&
仅可查看本人订单与本人通信内容；PID 作为系统内部匿名标识，默认不在普通页面展示；用户不直接接触真实身份与 PID 的映射关系。
\\

骑手端
&
查看分配订单；搜索配送任务；进入订单通信；处理配送异常；发起投诉申请。
&
仅可查看订单编号、订单状态、必要配送信息和通信窗口；不可见用户真实身份；不可见用户 PID；不可见用户备注和标签原文。
\\

平台服务器
&
用户认证；机器识别；PID 生成；订单绑定；Token 验证；消息转发；加密存储；条件溯源。
&
作为核心安全控制单元，维护必要的身份映射和订单状态；不向普通业务端暴露真实身份映射；敏感操作需要权限控制与审计记录。
\\

管理员端
&
管理员登录；订单检索；日志验证；处理溯源申请；查看审计结果。
&
默认只查看订单摘要、消息哈希、Merkle Root、日志验证结果和溯源申请状态；默认不展示通信明文、用户备注和标签原文；不能直接按 PID 任意溯源。
\\

审计日志系统
&
保存消息摘要；保存 Merkle Root 快照；记录管理员操作；支持日志完整性验证。
&
不作为明文消息查询入口；主要用于检测日志篡改，并为纠纷处理和管理员行为追踪提供可信依据。
\\

\bottomrule
\end{tabular}
\end{table}


从信任关系上看，FoodShield 采用平台可控匿名模型。平台服务器被视为系统中的基础可信实体，负责维护真实身份与 PID 之间的映射关系；用户和骑手彼此之间不完全信任，系统需要防止双方获取超过订单通信所需范围的敏感信息；管理员具有较高权限，但管理员行为也不能完全无约束，因此系统需要记录管理员的登录、查询、验证和溯源操作；外部攻击者不被信任，可能尝试伪造订单身份、重放 Token、篡改通信日志或越权访问管理员接口。

系统的主要信任边界包括以下几个方面：

\begin{enumerate}[label=\arabic*., leftmargin=2em]
\item \textbf{用户端与骑手端之间的边界}：用户与骑手只应围绕订单配送进行通信。骑手端不应获得用户真实身份和 PID，用户端也不应获得骑手端超出配送所需范围的敏感信息。
\item \textbf{普通业务端与管理员端之间的边界}：用户端和骑手端不能访问管理员接口；管理员端访问审计功能前必须完成身份认证。
\item \textbf{平台服务器与数据库之间的边界}：数据库保存订单、消息和审计数据，但通信明文不应直接裸露存储。消息内容应加密保存，消息摘要参与完整性验证。
\item \textbf{管理员权限与溯源能力之间的边界}：管理员不能任意根据 PID 发起溯源，而应基于订单号、投诉理由和审计需求触发条件溯源，并将操作写入审计日志。
\item \textbf{备注标签与审计数据之间的边界}：用户可为订单添加备注或标签以便历史订单管理和搜索，但管理员端默认不展示备注或标签的具体内容，避免审计权限扩大为隐私查看权限。
\item \textbf{日志数据与业务数据之间的边界}：日志审计系统主要保存消息摘要、Merkle Root 和操作记录，用于验证数据是否被篡改，而不是作为普通业务查询明文消息的入口。
\end{enumerate}

通过上述信任边界划分，FoodShield 将用户隐私保护、订单通信业务和平台管理能力进行分层隔离。不同参与方只能访问其业务所需的最小信息集合，从而降低身份暴露、越权访问、日志篡改和审计权限滥用风险。
